\documentclass[11pt,letterpaper]{article}
\usepackage[top=0.5in, bottom=0.38in, left=0.7in, right=0.7in]{geometry}
\usepackage{lmodern}
\usepackage{hyperref}
\hypersetup{colorlinks=false, urlcolor=black, pdfborder=0 0 0}
\setlength{\parindent}{0pt}
\setlength{\parskip}{0.55em}
\raggedright
\pagestyle{empty}

\begin{document}

\noindent
{\bfseries Ke-Cheng Hsu (Clio)}\\
Ph.D.\ Candidate, Department of Economics, HKUST\\
(852) 6840\,9753 \quad $|$ \quad \href{mailto:khsu@connect.ust.hk}{khsu@connect.ust.hk}\\
\href{https://cliohsu1997.github.io/}{cliohsu1997.github.io} \quad $|$ \quad \href{https://github.com/cliohsu1997}{github.com/cliohsu1997}

\vspace{0.3em}
\noindent
16 November 2026

\vspace{0.2em}
\noindent
Teacher Support Team\\
Registry Affairs Unit, HKU SPACE\\
Re: Part-time Teachers for Research and Data Storytelling (RF: 2026/27-PT-091-JDB)

\vspace{0.1em}
Dear Teacher Support Team,

I am writing to apply for the Part-time Teachers for Research and Data Storytelling post (RF: 2026/27-PT-091-JDB). I am a Ph.D.\ candidate in Economics at HKUST (expected June 2027), hold an M.Phil.\ in Economics (CGA 4.1/4.3), and am a Hong Kong permanent resident. I teach in English and Mandarin. The posting covers research methods, data analysis and interpretation, data visualisation, insight generation, and marketing / business decision-making---the same skills I use in doctoral research and in classroom teaching at HKUST.

I meet the degree and experience requirements. The post asks for a Master's degree in a related major (or specialist status) and at least three years of relevant teaching or industry experience. I completed an M.Phil.\ in Economics at HKUST and have taught continuously as a teaching assistant and lecturer since 2020---more than three years of relevant teaching on quantitative and research-methods topics. As lecturer for the MSc Mathematics Camp I received 4.53/5; as a TA for PhD and undergraduate economics my average evaluation is 4.72/5. I also served as Teaching Assistant Coordinator in the Business School (2025--2026), coaching others on clear materials and delivery.

I can teach the course topics with concrete research examples, not only textbook slides. Research methods and data analysis: I design empirical frameworks, define metrics, clean and link large panels, and estimate models with Python, R, and C++ on AWS/HPC. Data visualisation and interpretation: with Professors Kohei Kawaguchi and Kei Kawai I turned bookstore point-of-sale panels (tens of billions of records across 1{,}000+ stores) into weekly figures, regressions, and written summaries that show how demand uncertainty and inventory change cost and substitution---the storytelling step from raw data to a clear claim. Insight generation for marketing and business decisions: in my sole-authored job market paper on Taobao apparel return policy I framed a purchase--return trade-off, estimated structural responses under return-shipping insurance, and produced decision-ready comparisons (a moderate restocking-fee menu raises social welfare by about 2 million RMB, roughly a 5\% gain, and seller profit by about 1.2 million RMB relative to status-quo policies). That is exactly the loop of analysis, visualisation, and business recommendation the course title signals.

I am prepared to teach part-time for HKU SPACE with the same standards I use at HKUST: clear learning objectives, reproducible examples, and plain-language interpretation for learners who need research and data storytelling for professional decisions. My resume is attached. Thank you for your consideration. You can reach me at khsu@connect.ust.hk or +852 6840 9753.

\vspace{0.3em}
Sincerely,\\[0.25em]
Ke-Cheng Hsu (Clio)

\end{document}
